\documentclass[preview]{standalone}

\usepackage{amsmath}
\usepackage{amssymb}
\usepackage{stellar}
\usepackage{definitions}
\usepackage{bettelini}

\begin{document}

\id{power-series-calculus}
\genpage

\section{Termwise differentiation}

\begin{snippetproposition}{power-series-derivative-radius-proposition}{Radius of convergence of the differentiated series}
    Let
    \[
        \sum_{n=0}^{\infty}a_nx^n
    \]
    be a real power series with radius of convergence
    \(0<R\leq+\infty\). The termwise differentiated series
    \[
        \sum_{n=1}^{\infty}na_nx^{n-1}
    \]
    has the same radius of convergence \(R\). The two series may nevertheless
    behave differently at the boundary points \(x=\pm R\) when \(R<+\infty\).
\end{snippetproposition}

\begin{snippetproof}{power-series-derivative-radius-proposition-proof}{power-series-derivative-radius-proposition}{Radius of convergence of the differentiated series}
    Multiplication by \(x\) does not change the radius of convergence, so it
    suffices to compare \(\sum_na_nx^n\) with \(\sum_nna_nx^n\). By the
    Cauchy--Hadamard theorem,
    \[
        \frac1R=\limsup_{n\to\infty}\sqrt[n]{|a_n|}.
    \]
    Since \(\sqrt[n]{n}\to1\),
    \[
        \limsup_{n\to\infty}\sqrt[n]{n|a_n|}
        =\limsup_{n\to\infty}
        \left(\sqrt[n]{n}\sqrt[n]{|a_n|}\right)
        =\limsup_{n\to\infty}\sqrt[n]{|a_n|}.
    \]
    Thus both series have the same Cauchy--Hadamard radius.
\end{snippetproof}

\begin{snippettheorem}{power-series-termwise-differentiation-theorem}{Termwise differentiation of power series}
    Let \(f(x)=\sum_{n=0}^{\infty}a_nx^n\) have positive radius of
    convergence \(R\). Then \(f\) is infinitely differentiable on
    \((-R,R)\), and for every \(k\geq1\),
    \[
        f^{(k)}(x)
        =\sum_{n=k}^{\infty}
        n(n-1)\cdots(n-k+1)a_nx^{n-k}.
    \]
    Every derived series has the same radius of convergence \(R\).
\end{snippettheorem}

\begin{snippetproof}{power-series-termwise-differentiation-theorem-proof}{power-series-termwise-differentiation-theorem}{Termwise differentiation of power series}
    Fix \(0<r<R\). The original power series and its differentiated series
    converge uniformly on \([-r,r]\). The theorem on differentiation of a
    uniformly convergent sequence, applied to the partial sums, shows that the
    sum is differentiable there and that its derivative is the sum of the
    differentiated series. Since \(r<R\) was arbitrary, the conclusion holds
    throughout \((-R,R)\). Iterating the argument proves the formula for every
    \(k\).
\end{snippetproof}

\begin{snippetexample}{power-series-boundary-changes-after-differentiation-example}{Differentiation can change boundary convergence}
    The power series
    \[
        \sum_{n=1}^{\infty}\frac{x^n}{n}
    \]
    has radius \(R=1\) and converges on \([-1,1)\). Its differentiated series
    is
    \[
        \sum_{n=1}^{\infty}x^{n-1},
    \]
    which has the same radius but converges only on \((-1,1)\).
\end{snippetexample}

\begin{snippetexample}{power-series-successive-derivatives-boundary-example}{Successive derivatives and boundary convergence}
    The series
    \[
        \sum_{n=1}^{\infty}\frac{x^n}{n^2}
    \]
    has radius \(R=1\) and converges absolutely at both boundary points. Its
    first derivative,
    \[
        \sum_{n=1}^{\infty}\frac{x^{n-1}}n,
    \]
    converges on \([-1,1)\), while its second derivative,
    \[
        \sum_{n=2}^{\infty}\frac{n-1}{n}x^{n-2},
    \]
    converges only on \((-1,1)\).
\end{snippetexample}

\section{Recovery of coefficients}

\begin{snippetproposition}{power-series-coefficient-recovery-proposition}{Recovery of power-series coefficients}
    If
    \[
        f(x)=\sum_{n=0}^{\infty}a_nx^n
    \]
    has positive radius of convergence, then
    \[
        f^{(k)}(0)=k!a_k,
        \qquad
        a_k=\frac{f^{(k)}(0)}{k!}.
    \]
    Consequently, the values of a power-series function on any neighborhood
    of the center determine all its coefficients and therefore determine the
    function throughout its disk of convergence.
\end{snippetproposition}

\begin{snippet}{power-series-rigidity-interpretation}
    Power-series functions are rigid: two such representations that agree on
    a neighborhood cannot branch into different functions farther inside a
    common disk of convergence. This distinguishes analyticity from mere
    smoothness.
\end{snippet}

\begin{snippetexample}{smooth-flat-function-not-analytic-example}{A smooth function that is not analytic}
    Define
    \[
        f(x)=\begin{cases}
            e^{-1/x^2} & x\neq0,\\
            0 & x=0.
        \end{cases}
    \]
    This function is infinitely differentiable and
    \(f^{(k)}(0)=0\) for every \(k\geq0\). Its Maclaurin series is therefore
    identically zero, although \(f(x)>0\) for \(x\neq0\). Hence \(f\) is not
    analytic at the origin.
\end{snippetexample}

\section{Termwise integration}

\begin{snippettheorem}{power-series-termwise-integration-theorem}{Termwise integration of power series}
    Let \(f(x)=\sum_{n=0}^{\infty}a_nx^n\) have radius of convergence
    \(R>0\). For every \(|x|<R\),
    \[
        \integral[0][x][f(t)][t]
        =\sum_{n=0}^{\infty}\frac{a_n}{n+1}x^{n+1}.
    \]
    The integrated power series has the same radius of convergence \(R\).
\end{snippettheorem}

\begin{snippetproof}{power-series-termwise-integration-theorem-proof}{power-series-termwise-integration-theorem}{Termwise integration of power series}
    Fix \(r<R\). The power series converges uniformly on \([-r,r]\), so its
    partial sums may be integrated term by term on every interval contained in
    \([-r,r]\). This gives the displayed identity for \(|x|\leq r\). Since
    \(r<R\) is arbitrary, it holds for every \(|x|<R\). The coefficient
    ratio, or Cauchy--Hadamard theorem, shows that division by \(n+1\) does not
    change the radius.
\end{snippetproof}

\begin{snippetexample}{logarithm-from-integrated-geometric-series-example}{Logarithm from the geometric series}
    Integrating the geometric series on \(|x|<1\) gives
    \[
        \log(1-x)=-\sum_{n=1}^{\infty}\frac{x^n}{n}
    \]
    and, after replacing \(x\) by \(-x\),
    \[
        \log(1+x)=\sum_{n=1}^{\infty}(-1)^{n+1}\frac{x^n}{n}.
    \]
    The latter identity holds for \(-1<x<1\) and extends continuously to
    \(x=1\) by Abel's theorem.
\end{snippetexample}

\begin{snippetexample}{arctangent-from-integrated-geometric-series-example}{Arctangent from the geometric series}
    Since
    \[
        \frac1{1+x^2}=\sum_{n=0}^{\infty}(-1)^nx^{2n},
        \qquad |x|<1,
    \]
    termwise integration yields
    \[
        \arctan x=\sum_{n=0}^{\infty}(-1)^n
        \frac{x^{2n+1}}{2n+1},
        \qquad |x|<1.
    \]
    The series also converges at \(x=\pm1\), although its radius remains
    \(1\). The real function \(\arctan\) is analytic beyond this particular
    Maclaurin disk.
\end{snippetexample}

\end{document}
